\chapter{Related work}
\label{ch:related-work}

This thesis sits at the intersection of three lines of work: the radiance-field representations that
produced the Gaussian primitive (\Cref{sec:radiance-fields}), the volumetric Monte~Carlo methods
that make it physically renderable (\Cref{sec:volumetric-mc}), and the reference systems against
which the renderer is built and measured (\Cref{sec:reference-renderers}). The method it
re-implements and accelerates---\emph{Don't Splat Your Gaussians} (DSYG)---is reviewed in
\Cref{sec:dsyg}; the two ideas the contribution borrows from neighbouring work, decomposition
tracking and sort-free Gaussian rendering, are positioned alongside it.

\section{Radiance fields: from NeRF to Gaussian splatting}
\label{sec:radiance-fields}

Reconstructing a 3D scene from posed 2D photographs is the problem that produced the Gaussian
primitive. Neural Radiance Fields (NeRF)~\cite{NeRF} represent the scene \emph{implicitly}, as a
multilayer perceptron mapping a position and view direction to colour and density, rendered by
volumetric ray marching. NeRF synthesises photorealistic novel views from sparse input, and
follow-ups improve its anti-aliasing~\cite{MipNeRF} and few-shot behaviour~\cite{PixelNeRF}, but the
representation is a black box: it offers no explicit geometry to edit or relight, and both training
and rendering are slow.

3D Gaussian Splatting (3DGS)~\cite{3DGS} replaced the network with an \emph{explicit} set of
anisotropic 3D Gaussians, each carrying a position, covariance, opacity, and view-dependent colour.
Rendering rasterises the Gaussians---projecting them to screen space, sorting by depth, and
alpha-blending---which reaches real-time frame rates and high fidelity. The explicit representation
is directly editable, but it carries two costs relevant here. First, appearance is \emph{baked}: each
Gaussian stores a view-dependent colour fitted under the capture illumination, so the scene cannot be
relit. Second, correct alpha-blending requires a \emph{global depth sort} of the primitives, recomputed
per view; as the viewpoint moves the sort order changes discontinuously, producing the flickering and
popping characteristic of splatting. Both limitations point toward rendering the
same primitives by ray tracing instead of rasterisation.

\section{Ray-traced Gaussian primitives: Don't Splat Your Gaussians}
\label{sec:dsyg}

Condor et al.~\cite{DSYG} take the explicit Gaussians of 3DGS and render them as a
\emph{participating medium} by volumetric path tracing rather than rasterisation. Each Gaussian is
treated as a kernel of extinction whose contribution to the optical depth along a ray is evaluated
in closed form (\Cref{sec:kernel-mixture}), and the primitives are organised in a bounding volume
hierarchy. Because the renderer is physically based, the scene becomes \emph{relightable} and
supports multiple scattering and emission---recovering exactly what the baked, rasterised pipeline
gives up---while the closed-form optical depth keeps each primitive cheap. DSYG is the method this
thesis re-implements and accelerates.

The reference renders correctly but is not built for speed: it samples the scatter distance by
marching the ray boundary by boundary---maintaining the set of overlapping primitives and selecting
the next boundary by a running minimum---accumulating optical depth and solving for the scatter point
with a root-find (Newton or bisection) within each multi-primitive segment. That segment-by-segment
march and the per-segment root-find are precisely what \Cref{ch:architecture} removes, replacing them
with a single-trace collection pass and closed-form per-primitive free-flight sampling.

\paragraph{Concurrent ray-traced Gaussian methods.} Two concurrent works (2024) also ray-trace
Gaussian primitives, but toward radiance-field \emph{appearance} rather than physically-based
transport. 3D Gaussian Ray Tracing~\cite{3DGRT} replaces 3DGS rasterisation with a hardware-accelerated
BVH traversal over Gaussian particles, enabling secondary rays---reflections, refraction, distorted
cameras---while keeping 3DGS's alpha-composited, view-dependent colour model. EVER~\cite{EVER} renders
ellipsoids with \emph{exact} volumetric integration, removing the popping that approximate per-Gaussian
sorting introduces, but as an emissive--absorptive density field rather than a scattering medium with a
full light-transport path. This thesis instead accelerates DSYG's physically-based
\emph{multiple-scattering} path tracing of relightable participating media; its two contributions
(single-trace collection, argmin scatter sampling) target that transport problem, which neither
concurrent method addresses.

\section{Volumetric Monte Carlo and tracking}
\label{sec:volumetric-mc}

The estimators that make participating media renderable are surveyed by
Nov\'ak et al.~\cite{MonteCarlo}. For a heterogeneous medium the optical depth has no closed form, so
free paths are sampled without inverting it: \emph{delta (Woodcock) tracking}~\cite{Woodcock1965}
adds a fictitious medium to homogenise the extinction and accepts or rejects candidate collisions,
yielding unbiased free flights at the cost of many extinction evaluations. \emph{Decomposition} tracking~\cite{Kutz2017} reduces that cost by splitting the medium into a control
component---chosen so that free paths through it can be constructed in \emph{closed form}---and a
residual handled by tracking (analog in the variant this thesis builds on; weighted in its extension).

Decomposition tracking is the theoretical root of this thesis's scatter sampler. Its key property is
that a free flight through a sum of components is distributed as the \emph{minimum} of independent
free flights through each component taken separately. Because every Gaussian primitive here admits a
closed-form, invertible optical depth, each can draw its own free-flight distance directly, and the
minimum over the primitives a ray touches is the scatter event---no marching and no
root-finding (\Cref{ch:architecture}).

\section{Sort-free Gaussian rendering}
\label{sec:sort-free}

Removing the global sort from Gaussian rendering has also been pursued on the rasterisation side.
StochasticSplats~\cite{StochasticSplats2025} renders 3DGS scenes with stochastic rasterisation
(stochastic transparency), trading the depth sort for additional samples per pixel and thereby
avoiding the popping that sort-order changes introduce. The goal---avoiding an explicit
ordering of primitives---is the same one this thesis pursues, but in the complementary setting of a
physically based volumetric \emph{path tracer} rather than a rasteriser: here the ordering is sidestepped
by the decomposition-tracking argmin of \Cref{sec:adt} rather than by stochastic blending.

\section{Reference renderers and volume representations}
\label{sec:reference-renderers}

Mitsuba~3~\cite{Mitsuba3}, built on the Dr.Jit just-in-time compiler~\cite{DrJit}, is a research
renderer offering differentiable, retargetable (LLVM/CUDA) rendering from a Python front end; the
DSYG reference is implemented as a Mitsuba plugin. Its design priorities, however, leave performance
on the table for a bespoke production renderer: JIT kernel construction adds a per-process startup
cost---modest per run, but it compounds across the many short launches of a parameter sweep, where this
renderer's precompiled OptiX-IR pipeline starts markedly faster (\Cref{sec:results-startup} quantifies it); the
system is built for generality---differentiable, retargetable, inverse-rendering workflows---rather
than for a single, highly tuned forward renderer; and the generated CUDA offers limited room for the
fine-grained, primitive-specific optimisation this thesis applies. These gaps motivate a bespoke CUDA/OptiX implementation. Mitsuba nonetheless remains
indispensable as the \emph{validation reference}: \Cref{ch:validation} establishes correctness by
quantitative comparison against it.

Finally, the conventional representation for heterogeneous volumes is the sparse voxel grid, for
which OpenVDB~\cite{OpenVDB} and its GPU-friendly variant NanoVDB~\cite{NanoVDB} are the industry
standard. The Gaussian kernel mixture is a far more compact alternative---thousands of primitives
rather than millions of voxels---a difference \Cref{ch:results} quantifies in memory footprint.
